% !TEX program = xelatex
\documentclass[12pt]{article}

\input{Preamble}

\begin{document}

\section*{Аналитическое решение}

\subsection*{Постановка задачи №1}
В зал на $N$ мест заходят $N$ зрителей. Первый зритель --- Вася --- потерял билет и садится на любое случайно выбранное место. Каждый из последующих зрителей (со 2-го по $N$-го) действует по правилу: если его законное место свободно, он садится на него; если занято — выбирает любое свободное место случайным образом. Требуется найти вероятность того, что последний ($N$-й) зритель сядет на своё законное место.

\subsection*{Математическое решение №1}
Пусть $P(N)$ — вероятность того, что последний зритель сядет на своё место при общем количестве мест $N$. Вычислим её по формуле полной вероятности в зависимости от того, какое место изначально выбрал первый зритель:

\begin{enumerate}
    \item Вася сел на место №1 (вероятность $\frac{1}{N}$). Тогда все остальные спокойно садятся на свои места. Успех.
    \item Вася сел на место №$N$ (вероятность $\frac{1}{N}$). Тогда место последнего зрителя занято. Провал.
    \item Вася сел на место №$k$, где $1 < k < N$ (вероятность $\frac{1}{N}$ для каждого $k$). В этом случае зрители со 2-го по $(k-1)$-го садятся на свои места. Зритель $k$ обнаруживает своё место занятым и выбирает из оставшихся $(N - k + 1)$ мест (среди которых обязательно есть место №1 и место №$N$). Задача сводится к исходной, но для $(N - k + 1)$ человек.
\end{enumerate}

Получаем рекуррентное соотношение:
$$ P(N) = \frac{1}{N} \cdot 1 + \frac{1}{N} \cdot 0 + \sum_{k=2}^{N-1} \frac{1}{N} P(N - k + 1) $$

Сделаем замену переменной $j = N - k + 1$. При изменении $k$ от $2$ до $N-1$, индекс $j$ пробегает значения от $N-1$ до $2$. Перепишем сумму:
$$ P(N) = \frac{1}{N} \left( 1 + \sum_{j=2}^{N-1} P(j) \right) $$

Докажем по индукции, что $P(N) = \frac{1}{2}$.\\
База индукции для $N=2$: 
$$ P(2) = \frac{1}{2} (1 + 0) = \frac{1}{2} $$

Предположим, что $P(j) = \frac{1}{2}$ для всех $j < N$. Тогда шаг индукции для $P(N)$:
$$ P(N) = \frac{1}{N} \left( 1 + \sum_{j=2}^{N-1} \frac{1}{2} \right) $$
В сумме ровно $(N - 2)$ слагаемых. Раскрываем скобки:
$$ P(N) = \frac{1}{N} \left( 1 + \frac{N - 2}{2} \right) = \frac{1}{N} \left( \frac{2 + N - 2}{2} \right) = \frac{1}{N} \cdot \frac{N}{2} = \frac{1}{2}\ \square $$

\subsection*{Постановка задачи №2}
В зал на $N$ мест заходят $N$ зрителей. Первый зритель --- Вася --- потерял билет и садится на любое случайно выбранное место, \textbf{кроме своего собственного}. Каждый из последующих зрителей (со 2-го по $N$-го) действует по правилу: если его законное место свободно, он садится на него; если занято — выбирает любое свободное место случайным образом. Требуется найти вероятность того, что последний ($N$-й) зритель сядет на своё законное место.

\subsection*{Математическое решение №2}
Пусть $P(N)$ — вероятность того, что последний зритель сядет на своё место при общем количестве мест $N$ в новых условиях. Вычислим её по формуле полной вероятности в зависимости от того, какое место изначально выбрал первый зритель:

\begin{enumerate}
    \item Вася сел на место №1. По условию эта вероятность равна $0$, так как Вася принципиально не садится на своё место.
    \item Вася сел на место №$N$ (вероятность $\frac{1}{N-1}$). Тогда место последнего зрителя занято изначально. Успех невозможен (вероятность $0$).
    \item Вася сел на место №$k$, где $1 < k < N$ (вероятность $\frac{1}{N-1}$ для каждого $k$). В этом случае зрители со 2-го по $(k-1)$-го садятся на свои места. Зритель $k$ обнаруживает своё место занятым и выбирает из оставшихся $(N - k + 1)$ мест. Поскольку среди них есть место №1 (освобожденное Васей) и место №$N$, а выбор происходит вслепую, задача для оставшихся людей полностью сводится к классической, где вероятность успеха всегда равна $P_{\text{класс}} = \frac{1}{2}$.
\end{enumerate}

Получаем следующее соотношение:
$$ P(N) = \frac{1}{N-1} \cdot 0 + \frac{1}{N-1} \cdot 0 + \sum_{k=2}^{N-1} \frac{1}{N-1} P_{\text{класс}}(N - k + 1) $$

Подставим известное значение классической вероятности $P_{\text{класс}} = \frac{1}{2}$:
$$ P(N) = \frac{1}{N-1} \sum_{k=2}^{N-1} \frac{1}{2} $$

В сумме при изменении индекса $k$ от $2$ до $N-1$ содержится ровно $(N - 2)$ слагаемых. Раскрывая сумму, получаем итоговую формулу:
$$ P(N) = \frac{1}{N-1} \cdot \left( (N - 2) \cdot \frac{1}{2} \right) = \frac{N - 2}{2(N - 1)}\ \square $$

\end{document}